\section{Scoring rules and calibration}
\label{sec:calibration}

\subsection{What a probability forecast is scored on}
Write $\post=\Prob(Y=1\mid X=x)$ for the true spam probability and $q$ for a reported
forecast. Two forecasts can give the same spam label and assign very different probabilities
to an expensive mistake. For $N$ pairs $(q_i,y_i)$ the binary log-loss and Brier score are
\[
 -\frac{1}{N}\sum_{i=1}^N[y_i\log q_i+(1-y_i)\log(1-q_i)],
 \qquad \frac{1}{N}\sum_{i=1}^N(q_i-y_i)^2 .
\]
A loss is \emph{strictly proper} if reporting the true probability uniquely minimizes its
expected value. This concerns expectation over the outcome, even when an individual message
turns out to have been assigned a small probability.

\paragraph{Log-loss.}
For $0<\post,q<1$, the excess conditional loss over reporting $\post$ is
\[
 \post\log\frac{\post}{q}+(1-\post)\log\frac{1-\post}{1-q}.
\]
This is the Bernoulli Kullback--Leibler divergence. Applying $-\log u\geq 1-u$ to
$q/\post$ and $(1-q)/(1-\post)$ gives a lower bound of zero, with equality only at $q=\post$.
A forecast that assigns zero probability to a possible outcome has infinite expected loss.
For $\post=0$ the expected loss is $-\log(1-q)$, uniquely minimized at zero. For $\post=1$
it is $-\log q$, uniquely minimized at one.
\begin{assumptions}Binary outcomes, the natural logarithm, and the convention
$0\log 0=0$ for zero-weight terms. The proof uses the unclipped score. Numerical clipping
gives identical scores to forecasts clipped to the same value.\end{assumptions}

\paragraph{Brier score.}
Expanding the conditional expectation gives
\[
 \post(1-q)^2+(1-\post)q^2=\post(1-\post)+(q-\post)^2 .
\]
The excess loss is exactly $(q-\post)^2$, so the same unique minimizer applies on the
entire closed interval, including deterministic outcomes.
\begin{assumptions}$Y$ is binary and $q\in[0,1]$. The score uses the squared error for the
spam probability alone.\end{assumptions}

\subsection{Reliability, resolution and uncertainty}
Let $Q$ be a random forecast, $\eta(Q)=\E[Y\mid Q]$ the spam frequency conditional on that
forecast, and $\bar p=\E[Y]$. Calibration means $\eta(Q)=Q$ almost surely. Conditioning
on $Q$ makes the cross term in $(Q-Y)^2=[(Q-\eta)+(\eta-Y)]^2$ vanish.
The law of total variance then gives
\[
 \E[(Q-Y)^2]=
 \underbrace{\E[(Q-\eta(Q))^2]}_{\text{reliability penalty}}
 -\underbrace{\mathrm{Var}(\eta(Q))}_{\text{resolution}}
 +\underbrace{\bar p(1-\bar p)}_{\text{uncertainty}} .
\]
A constant forecast equal to the prevalence is calibrated but has zero resolution.
A forecast that separates groups with different spam frequencies gains resolution.
\begin{assumptions}The forecast and binary label refer to the same evaluation population.
Conditional means are taken at the forecast values actually issued.\end{assumptions}

For a reliability diagram, partition $[0,1]$ into disjoint bins. In each nonempty bin $b$,
plot its mean forecast $\bar q_b$ horizontally and its observed spam fraction $\bar y_b$
vertically. Points above the diagonal indicate underprediction. A count histogram shows how
many messages support each point. Equal-width bins use fixed probability intervals, while
equal-count bins aim for similar numbers of observations. The count-weighted expected
calibration error estimate is
\[
 \widehat{\mathrm{ECE}}=\sum_{b:n_b>0}\frac{n_b}{N}\,|\bar q_b-\bar y_b| .
\]
Under the empirical distribution, grouping identical forecasts gives the exact Brier identity
above. With wider bins, using $\bar q_b$ and $\bar y_b$ gives the exact bin-grouped decomposition
for forecasts replaced by their bin means. Within-bin variation can change the original
score, and opposite calibration errors can cancel in a reliability diagram.

\subsection{Identical labels, different bills}
Consider two groups of ten emails. The first contains two spam emails, the second eight.
Model A reports their observed frequencies $1/5$ and $4/5$. Model B sharpens these forecasts
by multiplying their log-odds by three, equivalently using temperature $T=1/3$.
The odds $1/4$ and $4$ become $1/64$ and $64$, so the forecasts become $1/65$ and $64/65$.
The transformation preserves the side of $0.5$ for every message. Both models assign
identical labels and correctly classify $16/20$ at that threshold.

With $(\cfp,\cfn,\ci)=(10,1,0.3)$ the inspection band is $(0.3,0.97)$. Both models deliver
the first group, at a cost of two missed spam emails. Model A inspects all ten emails
in the second group, giving total cost $2+10(0.3)=5$, or $0.25$ per email. Model B sends
the second group straight to spam. Its two false positives bring the total to $2+2(10)=22$,
or $1.10$ per email. The labels at $0.5$ hide the decisions that cross the upper
inspection threshold.
\begin{assumptions}The groups and labels are fixed for this realized-cost example.
Inspection always routes the message correctly and costs $0.3$ for either label.\end{assumptions}

\subsection{Experiments (C4)}
\label{sec:c4}
% Member 5. Every number below comes from python -m project4.experiments.task3_calibration
% (results/c4_scores.csv, results/c4_temperature.csv), using models/gaussian.py and models/logistic.py.

The example above used two hand-made groups of emails. We now measure the quality of the forecasts $q$ of
our two fitted classifiers on real and synthetic data, and check whether the same effect appears: equal
accuracy, different cost.

\paragraph{What we measure.}
On the test part of each split we compute the log-loss, the Brier score and $\widehat{\mathrm{ECE}}$ defined
above. For the log-loss we clip $q$ to $[10^{-12},\,1-10^{-12}]$, so that one confident mistake costs at most
$27.6$ instead of infinity. Reliability diagrams and ECE use $10$ bins, either equal-width
($[0,0.1),[0.1,0.2),\dots$) or equal-count (each holds about a tenth of the emails; equal values of $q$ are
always kept in the same bin). Next to these we report the accuracy of the labels ``spam iff $q\ge0.5$'' and,
for the Bayes rule with $\cfp=10$, $\cfn=1$, $\ci=0.3$, its realized loss per email and the fraction of emails
it inspects.

\paragraph{Setup.}
The \emph{synthetic} data are $4{,}601$ points (the size of Spambase) from two 2-D Gaussian classes with a
shared covariance, means $(0,0)$ and $(1.5,1)$, and $\pio=0.4$. Both models are correctly specified here,
so they should be calibrated. This is a check of the code, and the true posterior $\post$ gives a reference.
On \emph{Spambase} the features are transformed by $\log(1+x)$ and standardised, as described in the
experimental setup. The Gaussian model uses one shared covariance. Its covariance shrinkage and the L2 strength
of logistic regression are chosen on the validation part by \emph{log-loss}, because we want good
probabilities, not only good labels. All numbers are means over the 20 splits with standard errors
(Table~\ref{tab:c4-scores}).

\begin{table}[ht]
\centering\footnotesize\setlength{\tabcolsep}{3.5pt}
\input{\resultsdir/c4_scores.tex}
\caption{Forecast quality on the test part (mean $\pm$ standard error over 20 splits). Accuracy uses the
labels ``spam iff $q\ge0.5$''. Cost and Inspected are for the Bayes rule with $\cfp=10$, $\cfn=1$, $\ci=0.3$.
ECE uses 10 equal-width bins; equal-count bins give values within $0.003$. ``Bayes (true $p$)'' uses the
true posterior. ``+ Platt'' rows are re-calibrated on the validation part. Generated by
\texttt{experiments/task3\_calibration.py}.}
\label{tab:c4-scores}
\end{table}

\paragraph{Synthetic data.}
Both models score the same as the true posterior: log-loss $0.433$ against $0.432$, Brier score $0.139$ for
all three, and the same cost. Their reliability curves lie on the diagonal (Figure~\ref{fig:c4-rel}, top).
One detail matters when reading ECE: even the true posterior has $\widehat{\mathrm{ECE}}\approx0.03$. Each
split has only $919$ test points, so the observed spam fraction in each bin is noisy, and ECE adds up absolute
values of this noise. On data of this size, an ECE of about $0.03$ is the noise floor and does not mean the
model is miscalibrated.

\paragraph{Spambase.}
Logistic regression has better forecasts than the Gaussian model on every score:
\begin{itemize}[itemsep=1pt,topsep=2pt]
  \item log-loss $0.185$ against $0.224$;
  \item Brier score $0.050$ against $0.061$;
  \item ECE $0.021$ against $0.039$.
\end{itemize}
It also has the lower cost, $0.124$ against $0.142$ per email. The Gaussian model's reliability curve lies
above the diagonal almost everywhere, most of all for $q$ between $0.3$ and $0.6$ (Figure~\ref{fig:c4-rel},
bottom left). For example, emails it gives about $0.45$ are spam about $65\%$ of the time, so it
underpredicts spam in the middle range, which is where the inspection threshold $\thz=0.3$ lies. The reason
is the model's assumption that $p(x\mid y)$ is Gaussian. Spambase features are word counts with many zeros,
and even after $\log(1+x)$ they are far from Gaussian. Logistic regression models $\Prob(y\mid x)$ directly
and is trained by minimising the log-loss itself, so it is fitted to give good probabilities. Note also that
the two models differ by only $1.3$ points of accuracy but by $15\%$ in cost.

\begin{figure}[ht]
\centering
\includegraphics[width=0.92\linewidth]{c4_reliability.pdf}
\caption{Reliability diagrams on the test parts of all 20 splits pooled. Points on the dashed diagonal mean
the forecast matches the observed spam fraction; points above it mean the model underpredicts spam. The
histograms below each panel show how many forecasts fall in each equal-width bin (log scale).}
\label{fig:c4-rel}
\end{figure}

\paragraph{Same accuracy, different cost (temperature scaling).}
To repeat the example above on real data, we take the fitted logistic model on Spambase and change only
how confident it is:
\[
  q_T = \sigm\big(\logit q\,/\,T\big), \qquad T>0 .
\]
Dividing the log-odds by $T$ never changes their sign. So every email stays on the same side of $0.5$, and
the labels and the accuracy are \emph{exactly} the same for every $T$ (Table~\ref{tab:c4-temp},
Figure~\ref{fig:c4-temp}). $T<1$ pushes forecasts towards $0$ and $1$ (over-confident), and $T>1$ pulls
them towards $0.5$ (under-confident).

The decisions do change, because the Bayes rule compares $q_T$ with $0.3$ and $0.97$. At $T=0.4$ many
emails are pushed above $0.97$ and go straight to the spam folder: inspection falls from $20.6\%$ to $7.3\%$,
real emails sent to spam rise from $3.8$ to $10.5$ per test split (at a cost of $10$ each), and the cost
rises from $0.124$ to $0.164$ ($+33\%$). At $T=2.5$ the forecasts crowd around $0.5$, $46\%$ of all
emails are inspected, and the cost rises to $0.149$. Log-loss detects both changes ($0.289$ and $0.249$
against $0.185$). The Brier score hardly notices over-confidence ($0.053$ against $0.050$): a confident
mistake costs at most $1$ in squared error but $-\log q$ grows without bound. As in the example above, two
models with identical accuracy can have very different costs, and only a probability score can tell them
apart.

Close to $T=1$ the picture is less clear-cut. The cost is lowest at $T\approx1.2$ ($0.117$ against $0.124$ at
$T=1$), where a few more emails are inspected, and log-loss is lowest at $T\approx1.1$. So making the fitted
model slightly less confident would help a little. Platt scaling, below, finds a gain of the same size.

\begin{table}[ht]
\centering\small\setlength{\tabcolsep}{4pt}
\input{\resultsdir/c4_temperature.tex}
\caption{Temperature scaling of the logistic model on Spambase (test part, mean $\pm$ standard error over 20
splits). Accuracy is identical by construction; cost is for the Bayes rule with $\cfp=10$, $\cfn=1$,
$\ci=0.3$.}
\label{tab:c4-temp}
\end{table}

\begin{figure}[ht]
\centering
\includegraphics[width=\linewidth]{c4_temperature.pdf}
\caption{Accuracy, log-loss, realized loss and inspection rate of the logistic model on Spambase as its
temperature changes (bands: $\pm$ one standard error over 20 splits). Accuracy is flat. Log-loss and cost
rise clearly once $T$ is far from $1$; both are lowest slightly above $T=1$.}
\label{fig:c4-temp}
\end{figure}

\paragraph{Re-calibrating on validation data (Platt scaling).}
If a model's forecasts are off, we can repair them without refitting the model. We fit
\[
  q' = \sigm\big(A\,\logit q + B\big)
\]
on the \emph{validation} part. This is a logistic regression with the single feature $\logit q$ plus an
intercept, so we fit it with the same Newton/IRLS update \eqref{eq:t3-irls}. $A$ plays the role of $1/T$,
$B$ shifts the forecasts up or down, and $A=1$, $B=0$ leaves $q$ unchanged. For the Gaussian model on
Spambase this lowers ECE from $0.039$ to $0.024$, log-loss from $0.224$ to $0.208$, and the cost from $0.142$
to $0.132$. The logistic model was already close to calibrated: its log-loss hardly changes, and its cost
falls slightly, from $0.124$ to $0.118$. On the synthetic data, where both models are already calibrated,
Platt scaling changes nothing. Because $B$ moves the point where $q'=0.5$, accuracy can change a little.

\begin{assumptions}
Test emails come from the same distribution as the training and validation emails (no prior shift; that
is Section~\ref{sec:prior-shift}). ECE depends on the number of bins and has a positive bias on small
test sets. Platt scaling assumes that a straight-line correction of the log-odds is enough. It is fitted on
the same validation part that chose the hyperparameters, so its gain may be slightly optimistic; the test
part is never used to choose anything.
\end{assumptions}

\begin{sloppypar}
\emph{Code and checks.} \texttt{calibration.py} implements every score and both re-scalings.
The unit tests in \texttt{tests/test\_calibration.py} check the following:
\end{sloppypar}
\begin{itemize}[itemsep=1pt,topsep=2pt]
  \item the log-loss on the worked example of the lecture notes;
  \item log-loss and Brier score against scikit-learn;
  \item the reliability bins on a small example worked by hand, and that equal values of $q$ are never
        split between bins;
  \item that ECE is near $0$ for data that are calibrated by construction, and large after over-confident
        scaling;
  \item that on calibrated data no temperature beats $T=1$ for either score, as properness predicts;
  \item that Platt scaling recovers a known $(A,B)$ and matches scikit-learn's logistic regression.
\end{itemize}
